\begin{theorem}[I9 candidate: continuity in degree one for arbitrary complete targets]
Let $A$ be a regular noetherian local ring of finite dimension, let
$K$ be its residual resolution, and put
$B=\prod_{\mathfrak p}\widehat{A_{\mathfrak p}}$ with its product
adic topology. Let $Q$ be a separated complete graded topological
$A$-algebra with jointly continuous multiplication, a continuous
degree-zero homomorphism from $B$, and $T\in Q^1$ satisfying the
immediate-specialization incidence relations. Then there is a unique
continuous additive map
\[
 \Phi_1:\operatorname{End}_A^1(K)\longrightarrow Q^1
\]
which sends $b\theta c$ to $bTc$ for $b,c\in B$.
It respects the two $B$-actions. No open-left-ideal hypothesis is used.
This statement concerns degree one only; it does not establish the
full presentation in higher degrees.
\end{theorem}
\begin{proof}
Write $T_{\mathfrak p}=\widehat{A_{\mathfrak p}}$. For an adjacent
pair $\mathfrak p\subset\mathfrak q$, Proposition I4 says that
\[
 T_{\mathfrak p}\longrightarrow
 \operatorname{Hom}_A(E_{\mathfrak p},E_{\mathfrak q}),
 \qquad c\longmapsto\delta_{\mathfrak p\mathfrak q}c
                                                        \tag{I9.1}
\]
is injective with dense image. Its induced topology has a basis
\[
 (\mathfrak qA_{\mathfrak q})^N A_{\mathfrak q}
       +(\mathfrak pA_{\mathfrak p})^nT_{\mathfrak p}.
                                                        \tag{I9.2}
\]
The map $c\mapsto e_{\mathfrak q}Tc$ is continuous for this mixed
topology: split (I9.2) into its two terms, use the left completed
$A_{\mathfrak q}$-action for the first, and the prescribed right
$T_{\mathfrak p}$-action for the second. It therefore extends to
the completion in (I9.1). Denote this continuous map by
$\phi_{\mathfrak q\mathfrak p}$.

Fix a source height $h$. Let $D_h$ be the subgroup of
$\operatorname{Hom}_A(K^h,K^{h+1})$ consisting of finite matrices
whose adjacent entries have the form (I9.1). It is dense for the
finite topology: truncate the finitely many source and target
summands needed for prescribed inputs, then apply I4 to each of
the finitely many entries. Injectivity in (I9.1) defines an
unambiguous additive map
\[
 \phi:D_h\longrightarrow Q^1,
 \qquad (\delta_{\mathfrak p\mathfrak q}c_{\mathfrak q\mathfrak p})
       \longmapsto
       \sum_{\mathfrak p,\mathfrak q}
                    e_{\mathfrak q}Tc_{\mathfrak q\mathfrak p}.
                                                        \tag{I9.3}
\]
We prove that this map is continuous for the subspace finite
topology. All limits below involve only finitely many adjacent
entries at once, so their images converge by the individual-arrow
continuity just established.

\emph{Uniform control outside finitely many source primes.}
Let $W$ be a closed zero-neighborhood in $Q^1$. Joint continuity
gives product-adic neighborhoods $V_L,V_R\subset B$ such that
\[
                         V_L T V_R\subset W.              \tag{I9.4}
\]
Choose finite sets of primes and powers defining basic product
neighborhoods inside $V_L,V_R$. Let $P_0$ be the finite set of
height-$h$ primes constrained by $V_R$. We claim that every finite
matrix in $D_h$ supported on source primes outside $P_0$ has image
in $W$.

Indeed, fix its finite source set $P$ and finite target set $Q_1$.
The semilocal ring
\[
 A_P=(A\setminus\bigcup_{\mathfrak p\in P}\mathfrak p)^{-1}A
\]
has precisely the incomparable primes $\mathfrak pA_P$, for
$\mathfrak p\in P$, as its maximal ideals. Chinese remaindering
therefore approximates, for each target $\mathfrak q$, its entire
finite row of coefficients by a single element $a_{\mathfrak q}$
of $A_P$, to any prescribed $\mathfrak p$-adic orders at the
sources. Values at nonincident pairs may be prescribed arbitrarily.
Clear all the finitely many denominators by one $u$ outside every
prime in $P$, so that $u a_{\mathfrak q}\in A$.

Choose additionally $v$ outside every prime in $P$, lying in each
of the finitely many target-prime powers prescribed by $V_L$.
Such a $v$ exists by prime avoidance: every constrained target
prime has height $h+1$, so its power is contained in none of the
height-$h$ source primes. A product of choices works for their
intersection. Set
\[
 b_{\mathfrak q}=vu a_{\mathfrak q}\quad(\mathfrak q\in Q_1),
 \qquad c_{\mathfrak p}=(vu)^{-1}\quad(\mathfrak p\in P),
\]
and set all other components equal to zero. Then $b\in V_L$ and
$c\in V_R$, because $P$ avoids $P_0$. The adjacent coefficients
of $b\theta c$ are $a_{\mathfrak q}$ at the specified sources,
and hence approximate the original finite matrix. The corresponding
values $bTc$ lie in $W$ by (I9.4). Individual-arrow continuity and
closedness of $W$ prove the claim. This bound is independent of the
number of source or target primes in the finite matrix.

\emph{Continuity on a fixed source column.}
Fix $\mathfrak p$ of height $h$, and write
$D_{\mathfrak p}\subset D_h$ for its finite-support column.
We show that the restriction of (I9.3) to this column is continuous.
Choose closed neighborhoods sufficiently small that the finite sums
used below lie in a prescribed neighborhood. In particular choose
$V_L,V_R$ so that a sum of at most $h$ elements of $V_LTV_R$ is
small. If $h=0$, this remainder sum is empty.

Choose parameters $x_1,\ldots,x_h\in A$ for
$A_{\mathfrak p}$ and $n$ large enough that a vector supported
at $\mathfrak p$ with coefficient in
$(\mathfrak pA_{\mathfrak p})^nT_{\mathfrak p}$ belongs to $V_R$.
Take the standard cyclic generator $e\in E_{\mathfrak p}$ with
\[
 I=\operatorname{Ann}_A(e),\qquad
 I A_{\mathfrak p}=(x_1^n,\ldots,x_h^n)A_{\mathfrak p}.
                                                        \tag{I9.5}
\]
For $h=0$, take $e=1\in\operatorname{Frac}(A)$ and $I=0$.
The ideal $I$ is $\mathfrak p$-primary. The finite module
\[
                         M=\operatorname{Ext}_A^h(A/I,A)
\]
identifies with the kernel of the first differential in
$\operatorname{Hom}_A(A/I,K)$, and thus embeds into
$E_{\mathfrak p}[I]$. Choose finitely many generators
$m_1,\ldots,m_r$ of $M$. The map
$T_{\mathfrak p}/IT_{\mathfrak p}\to E_{\mathfrak p}[I]$,
$c\mapsto ce$, is an isomorphism: the parameter-power quotient
in (I9.5) is an Artinian complete intersection. Hence write
$m_i=c_i e$ with fixed $c_i\in T_{\mathfrak p}$.

For an adjacent $\mathfrak q$, localization of this cohomology
description gives
\[
 M_{\mathfrak q}
   =\ker\bigl(E_{\mathfrak p}[I]
                \xrightarrow{\delta_{\mathfrak p\mathfrak q}}
                         E_{\mathfrak q}[I]\bigr).        \tag{I9.6}
\]
There are no other height-$h$ primes over the $\mathfrak p$-primary
ideal, and after localization at $\mathfrak q$ there is only the
one possible height-$h+1$ target. This proves (I9.6) directly.

Continuity of the map
\[
 (b_1,\ldots,b_r)\longmapsto
          \sum_i b_i T e_{\mathfrak p}c_i
\]
from the finite product of target-diagonal rings gives a finite
set $Q_0$ of target primes such that all tuples supported outside
$Q_0$ have small image. Enlarge $Q_0$ to include all target
coordinates constrained by $V_L$.

Suppose $f\in D_{\mathfrak p}$ vanishes on $e$ and is supported
outside $Q_0$. Write its finite entries as
$\delta_{\mathfrak p\mathfrak q}c_{\mathfrak q}$.
By (I9.6), for every such $\mathfrak q$ there are
$b_{\mathfrak q i}\in A_{\mathfrak q}$ such that
\[
 c_{\mathfrak q}
    =\sum_i b_{\mathfrak q i}c_i+r_{\mathfrak q},
 \qquad r_{\mathfrak q}\in IT_{\mathfrak p}.              \tag{I9.7}
\]
The first terms have small total image by the definition of $Q_0$.
For the remainders write
\[
 r_{\mathfrak q}=\sum_{\ell=1}^h x_\ell^n a_{\mathfrak q\ell},
 \qquad a_{\mathfrak q\ell}\in T_{\mathfrak p}.
\]
Approximate these finitely many coefficients by elements of
$A_{\mathfrak p}$ and clear all denominators by one $s\notin
\mathfrak p$. The resulting remainder column is a sum of at most
$h$ words
\[
             b_\ell\theta c_\ell,\qquad
 (c_\ell)_{\mathfrak p}=x_\ell^n/s,
\]
where $b_\ell$ is supported on the finitely many target primes
outside $Q_0$ and has coefficients in $A$. Thus $b_\ell\in V_L$
and $c_\ell\in V_R$. Their images have small sum. Passing to the
limit in the finitely many adjacent entries preserves the bound,
by the chosen closed neighborhoods.

Finally, the finitely many target components in $Q_0$ are controlled
by the already proved individual-arrow continuity. Require $f$ to
vanish on $e$ and on the finitely many additional inputs supplied
by those components. Splitting $f$ into its $Q_0$ part and its
complement proves continuity on the whole fixed source column.

\emph{Assembly and completion.}
For a given target neighborhood, the uniform-tail argument leaves
only finitely many source primes $P_0$. Apply the fixed-column
argument to each of them. A finite union of their finite input
sets is a valid finite-topology neighborhood on $D_h$. Thus
(I9.3) is continuous. There are only finitely many heights $h$,
so the maps combine into a continuous map on the dense finite-arrow
subgroup of $\operatorname{End}_A^1(K)$. An additive continuous map
is uniformly continuous; completeness and separation of $Q^1$
give a unique extension to the whole degree-one endomorphism group.

The formulas agree with both $B$-actions first on finite entries,
then on all entries by continuity. Finite diagonal truncations of
$b\theta c$ converge to it in the finite topology and their
images converge to $bTc$ by joint continuity in $Q$. Therefore
$\Phi_1(b\theta c)=bTc$. Density proves uniqueness.
\end{proof}

\paragraph{Status and remaining gate.}
This is a new candidate proof by the delegated BCE agent, awaiting a
separate AI audit. It uses the proved I4 individual-arrow theorem and
the ordinary finite-module computation displayed above. It claims no
solution in higher degrees. Even after acceptance of degree-one
continuity, the construction of the higher-degree map and the full
closed relation ideal in an arbitrary complete target remain to be
proved.
